% ============================================================================
% Chapitre 0 — Introduction générale
% ============================================================================
\chapter{Introduction générale}
\markboth{Introduction générale}{Introduction générale}

\section*{Introduction}

L'essor fulgurant des architectures \textit{cloud-native} au cours de la dernière
décennie a profondément transformé la manière dont les organisations déploient et
gèrent leurs systèmes informatiques. Les infrastructures modernes, composées de
dizaines voire de centaines de serveurs, de conteneurs et de services distribués,
génèrent des flux de données opérationnelles d'une densité sans précédent. Dans ce
contexte, la supervision des systèmes --- autrefois déléguée à un administrateur
réseau scrutant quelques courbes --- est devenue une discipline à part entière,
au cœur de la résilience organisationnelle.

Pourtant, les outils actuels de surveillance, aussi sophistiqués soient-ils,
partagent une limite fondamentale : ils savent \textit{détecter} qu'un problème
existe, mais peinent à expliquer \textit{pourquoi} il est apparu. Face à une alerte
« CPU à 98~\% », l'ingénieur de garde doit encore mener une investigation manuelle
qui peut durer de trente minutes à plusieurs heures, consultant successivement les
journaux système, l'historique des déploiements, les requêtes base de données et les
métriques réseau. Ce gouffre entre \textit{détection} et \textit{compréhension}
constitue précisément le problème que ce projet cherche à résoudre.

\section{Contexte et motivation}

Les statistiques sectorielles confirment l'ampleur du défi. Selon le rapport
\textit{State of Digital Operations} de PagerDuty (2023) \cite{pagerduty2023},
le coût moyen d'une heure d'indisponibilité s'élève à plus de 300~000 dollars
pour les grandes entreprises, tandis que le temps moyen de résolution (\textsc{mttr})
dépasse quarante-cinq minutes dans les organisations ne disposant pas de pratiques
\textsc{aiops} matures.

Pour les petites équipes --- une startup avec deux développeurs gérant un serveur
\textsc{vps}, une agence web administrant une dizaine de machines --- la situation
est encore plus critique. Ces équipes n'ont ni la ressource humaine d'un ingénieur
\textit{DevOps} dédié, ni le budget pour des solutions commerciales à plusieurs
milliers d'euros par mois. Elles ont besoin d'un outil qui réponde simplement à
la question : « \textit{Mon serveur est-il en bonne santé ? Sinon, que dois-je
faire ?} »

C'est dans ce contexte que s'inscrit le projet \textbf{\caimp{}} ---
\textit{Centralized AI Monitoring Platform} --- une plateforme auto-hébergée
et \textit{open source} qui positionne l'intelligence artificielle non pas comme
un gadget d'affichage, mais comme le cœur du processus de diagnostic.

\section{Objectifs du projet}

Les objectifs du projet \caimp{} ont été définis de manière itérative en tenant
compte des lacunes identifiées dans les solutions existantes :

\begin{enumerate}[label=\textbf{O\arabic*.}, leftmargin=2.2cm]
  \item \textbf{Tableau de bord orienté réponses} : concevoir une interface dont
    la page d'accueil présente immédiatement le statut global, les problèmes
    potentiels et les actions recommandées --- sans nécessiter d'expertise
    en monitoring de la part de l'utilisateur.

  \item \textbf{Corrélation automatique des changements} : développer un moteur
    capable d'associer chaque anomalie détectée aux modifications récentes de
    l'environnement (déploiements, redémarrages, mises à jour de paquets,
    connexions \textsc{ssh}) afin d'automatiser la recherche de la cause racine.

  \item \textbf{Explication en langage naturel} : intégrer un modèle de langage
    local (sans appel à un \textsc{api} cloud) pour produire une explication
    compréhensible par un développeur sans formation \textit{DevOps}.

  \item \textbf{Prévision des saturations} : implémenter un moteur de prévision
    basé sur le lissage exponentiel de Holt-Winters pour anticiper les dépassements
    de seuils dans les 24 heures à venir.

  \item \textbf{Architecture microservices robuste} : concevoir une architecture
    découplée, basée sur un bus d'événements (\textsc{nats} JetStream), permettant
    la scalabilité horizontale et l'isolation multi-tenant des données.

  \item \textbf{Sécurité de bout en bout} : implémenter l'authentification mutuelle
    \textsc{tls} pour les agents, la gestion de jetons \textsc{jwt}, et l'isolation
    des données par organisation via des politiques \textsc{rls} en base de données.

  \item \textbf{Déploiement accessible} : rendre l'installation complète réalisable
    en moins de dix minutes via Docker Compose, sur n'importe quelle machine disposant
    de 6~Go de \textsc{ram}.
\end{enumerate}

\section{Méthodologie de travail}

Le développement de \caimp{} a suivi une approche \textit{agile} adaptée au
contexte académique. L'équipe de trois personnes s'est réunie quotidiennement pour
des sessions de synchronisation courtes (\textit{stand-ups}), et a organisé son
travail en sprints hebdomadaires d'une durée de cinq jours.

Le dépôt de code a été hébergé sur GitHub, avec une branche principale protégée
(\texttt{main}) et une branche de développement (\texttt{dev}) depuis laquelle
les fonctionnalités ont été développées et testées avant intégration. Les messages
de \textit{commit} ont suivi la convention \textit{Conventional Commits}
(\texttt{feat:}, \texttt{fix:}, \texttt{chore:}, \texttt{docs:}), facilitant la
génération automatique du journal des modifications.

La répartition des responsabilités a été organisée par domaine d'expertise :
le \textit{backend} (services Python, Go, base de données) a mobilisé les trois
membres en rotation, le \textit{frontend} React a été principalement pris en
charge par un membre spécialisé, et l'intégration des modèles \textsc{llm} et
du moteur de corrélation a constitué un chantier partagé en binôme.

\section{Plan du rapport}

Ce rapport est structuré en dix chapitres, organisés pour guider le lecteur
depuis le contexte général jusqu'aux perspectives d'évolution :

\begin{itemize}
  \item Le \textbf{chapitre~1} présente le contexte général du projet, les enjeux
    de l'observabilité moderne et le domaine de l'\textsc{aiops}.
  \item Le \textbf{chapitre~2} analyse l'état de l'art des solutions de supervision
    existantes et positionne \caimp{} par rapport à elles.
  \item Le \textbf{chapitre~3} spécifie les besoins fonctionnels et non fonctionnels,
    et présente les diagrammes UML d'analyse.
  \item Le \textbf{chapitre~4} décrit en détail l'architecture de la plateforme,
    le modèle de données et les flux de traitement.
  \item Le \textbf{chapitre~5} justifie les choix technologiques.
  \item Le \textbf{chapitre~6} présente les détails d'implémentation des modules
    clés, avec extraits de code commentés.
  \item Le \textbf{chapitre~7} documente la stratégie de tests et les résultats
    de validation.
  \item Le \textbf{chapitre~8} décrit les procédures de déploiement et de mise
    en production.
  \item Le \textbf{chapitre~9} discute les résultats obtenus et les perspectives
    d'évolution.
  \item Le \textbf{chapitre~10} conclut le rapport et dresse un bilan personnel
    de l'expérience.
\end{itemize}

\section*{Conclusion}

Ce projet représente l'aboutissement d'un mois de travail intense, pendant lequel
une plateforme complète de supervision d'infrastructure a été conçue, développée,
testée et documentée. Les chapitres qui suivent témoignent de la rigueur
scientifique et technique que nous avons cherché à maintenir tout au long de cette
démarche. Le chapitre suivant pose les bases contextuelles indispensables à la
compréhension de l'ensemble du travail réalisé.
